\documentclass[10pt,twoside]{article}
\usepackage[a4paper,inner=24mm,outer=17mm,top=18mm,bottom=18mm]{geometry}
\usepackage{multicol}
\usepackage{xcolor}
\usepackage{lmodern}
\usepackage{fancyhdr}
\usepackage{titlesec}
\usepackage[hidelinks]{hyperref}
\definecolor{zineink}{HTML}{252525}
\definecolor{zineorange}{HTML}{C85C32}
\pagestyle{fancy}
\fancyhf{}
\fancyhead[LE,RO]{\sffamily\small\color{zineorange}FIELD GUIDE}
\fancyhead[RE,LO]{\sffamily\small\color{zineink}ISSUE 01}
\fancyfoot[C]{\sffamily\small\thepage}
\renewcommand{\headrulewidth}{0pt}
\titleformat{\section}{\sffamily\large\bfseries\color{zineorange}}{}{0pt}{}
\setlength{\columnsep}{8mm}
\setlength{\parindent}{0pt}
\setlength{\parskip}{4pt}
\begin{document}
\begin{center}
  {\sffamily\fontsize{30}{32}\selectfont\bfseries FIELD GUIDE}\par
  {\sffamily\large\color{zineorange}A minimal two-column zine for small adventures}\par
  \smallskip\rule{0.9\linewidth}{1pt}
\end{center}
\begin{multicols}{2}
\section*{The quiet path}
The best route is rarely the straightest. A path with one useful landmark,
one place to pause, and one question worth asking is enough for a first
expedition.

\section*{Three tools}
\textbf{A pencil.} Mark what changes when you return.\par
\textbf{A timer.} Notice how long a choice takes.\par
\textbf{A friend.} Compare what each of you saw.

\section*{Encounter table}
\begin{tabular}{p{0.18\columnwidth}p{0.48\columnwidth}}
1 & A gate left open\\
2 & A useful shadow\\
3 & A note under a stone\\
4 & A path that ends early
\end{tabular}

\section*{Make it yours}
Replace this page with maps, tables, character notes, or a short adventure.
The inner and outer margins are deliberately explicit so the document can be
imposed as a folded booklet without losing the binding edge.
\end{multicols}
\end{document}
